% =====================================================================
%  CUERPO DEL INFORME TÉCNICO  ---  MÁXIMO 20 PÁGINAS
% =====================================================================
%  CÓMO USAR ESTE ARCHIVO
%  Cada subsección trae un bloque de comentarios con:
%     [CONTENIDO]  -> los puntos que el párrafo DEBE cubrir (mapeados a la rúbrica)
%     [ESTRUCTURA] -> cómo encadenar los párrafos
%     [EXTENSIÓN]  -> presupuesto de palabras para no pasar las 20 páginas
%  Escribe tu texto donde dice ESCRIBIR AQUÍ y borra el comentario cuando termines.
%
%  PRESUPUESTO DE PÁGINAS (medido con esta plantilla: ~360 palabras por plana)
%     Figuras, tablas y títulos ya insertados .. ~9 planas (verificado al compilar)
%     Texto disponible ......................... ~11 planas = ~4.000 palabras
%     TOTAL objetivo ........................... 20 planas (límite duro de la rúbrica)
%  Los [EXTENSIÓN] de cada subsección suman aproximadamente ese total. Compila
%  cada cierto tiempo y revisa dónde cae el título "4. Referencias bibliográficas":
%  debe quedar en la plana 25 del PDF o antes (portada + abstract + 3 índices = 5
%  planas previas al cuerpo). Si te pasas, el orden de recorte es:
%     1) fusionar 2.4.1 y 2.4.2 en un solo párrafo;
%     2) eliminar la Figura 3-2;
%     3) achicar la Figura 2-2 a width=0.75\textwidth.
%
%  REGLAS DE REDACCIÓN (rúbrica E)
%   - Tercera persona con sujeto explícito: "el practicante desarrolló",
%     "el proyecto contempló", "DTS definió". NO uses voz pasiva ni "se hizo".
%   - Un punto aparte por idea central. Nada de párrafos de una sola oración.
%   - Cursiva SOLO para palabras en otro idioma: \textit{Store \& Forward},
%     \textit{link budget}, \textit{payload}, \textit{watchdog}.
%   - Toda figura y tabla debe estar mencionada en el texto antes de aparecer.
% =====================================================================


% =====================================================================
\section{Identificación y Contextualización}
% =====================================================================
% [EXTENSIÓN TOTAL DE LA SECCIÓN] ~4,5 planas (1.020 palabras + Figura 1-1 + Tabla 1-1).

% ---------------------------------------------------------------------
\subsection{Contextualización de la empresa}
% ---------------------------------------------------------------------
% [!! LÍMITE DURO: UNA PLANA. La rúbrica lo dice explícitamente. Apunta a 300 palabras.]
%
% [CONTENIDO]  (rúbrica B / Identificación, criterio 1)
%   - Identidad legal: DTS SpA, constituida el 11 de enero de 1991, transformada
%     a sociedad por acciones en noviembre de 2020, empresa relacionada de ENAER.
%   - Giro y rubro: industria electrónica; defensa, seguridad y aeroespacial.
%   - TAMAÑO EN RELACIÓN AL RUBRO (la rúbrica lo pide explícitamente):
%     dotación en el rango 50-250 trabajadores, clasificada como gran empresa
%     según el criterio de ventas anuales del Estatuto Pyme (Ley 20.416).
%     >>> Este es el criterio que hay que citar: la clasificación chilena es por
%         ventas en UF, no por dotación. Cítalo con (Ministerio de Economía,
%         Fomento y Turismo, 2014) -- es la misma fuente que usaste en Práctica I.
%   - Áreas de negocio: nómbralas en una sola oración, sin viñetas, para ahorrar
%     espacio. No copies el folleto completo.
%   - Certificaciones que dan contexto de rigor técnico: ISO 9001:2015,
%     FAA N1DSY390C (FAR 145), DGAC N268 (DAN 145), laboratorio de calibraciones
%     acreditado NCh-ISO/IEC 17025:2017.
%   - Ubicación: casa matriz en Macul, oficinas técnicas en Iquique, Antofagasta,
%     Talcahuano y Punta Arenas, oficina comercial en Viña del Mar.
%   - CIERRE OBLIGATORIO: el área I+D+i y el hecho de que el desarrollo espacial
%     era una línea SIN trayectoria previa. Este cierre es el que conecta con 1.3
%     y con el problema; no lo dejes fuera.
%
% [ESTRUCTURA] 3 párrafos:
%   P1  Quién es DTS: identidad, giro, trayectoria, relación con ENAER, tamaño
%       y su clasificación en el rubro (con cita).
%   P2  Qué hace: áreas de negocio, certificaciones y cobertura geográfica.
%   P3  Dónde se insertó el practicante: departamento I+D+i, su rol dentro de la
%       empresa y la ausencia de desarrollo espacial previo.
%
% [EJEMPLO DE ARRANQUE DE P1]
%   "Desarrollo de Tecnologías y Sistemas SpA (DTS) es una empresa chilena
%    constituida en 1991 y relacionada a la Empresa Nacional de Aeronáutica de
%    Chile (ENAER), cuyo giro principal corresponde a la industria electrónica
%    aplicada a los mercados de defensa, seguridad y aeroespacial. Con una
%    dotación cercana a los [N] trabajadores, la empresa se clasifica como gran
%    empresa según el criterio de ventas anuales establecido en el Estatuto Pyme
%    (Ministerio de Economía, Fomento y Turismo, 2014), condición poco frecuente
%    en un rubro nacional dominado por proveedores de menor escala."

% ESCRIBIR AQUÍ (máx. 1 plana)




% ---------------------------------------------------------------------
\subsection{Descripción del proyecto y rol del practicante}
% ---------------------------------------------------------------------
% [EXTENSIÓN] ~300 palabras + Figura 1-1 + Tabla 1-1.
%
% [CONTENIDO]  (rúbrica B / Identificación, criterio 2: objetivo general, plazos,
%               alcance, descripción técnica, procesos y agentes involucrados)
%   - Nombre del proyecto: BIFROST (propuesto por el propio practicante -- vale
%     la pena decirlo, aporta a "aplicación creativa" de la evaluación global).
%   - OBJETIVO GENERAL DEL PROYECTO (declararlo de forma explícita y en una sola
%     oración, porque la rúbrica lo busca literalmente):
%       "diseñar y validar, a nivel de arquitectura, una carga útil de
%        comunicaciones tipo Store & Forward para un satélite CubeSat 12U, capaz
%        de recibir mensajes autónomos desde estaciones remotas terrestres en
%        banda UHF, almacenarlos a bordo y descargarlos a estaciones terrenas."
%   - Descripción técnica de alto nivel: receptor UHF 401-403 MHz, memoria no
%     volátil a bordo, downlink UHF ~466 MHz por antena propia de la carga útil,
%     bus CAN hacia el OBC, órbita SSO 650 km / 96 grados de inclinación.
%   - PLAZOS: práctica del 2 de enero al 28 de febrero de 2026, jornada completa.
%     Fases cubiertas: Pre-Phase A, Phase A (cierre en SRR, 30-01-2026) y
%     Phase B (cierre en PDR, 27-02-2026).
%   - ALCANCE Y NO-ALCANCE (decirlo explícito suma): queda FUERA el diseño
%     detallado (CDR), la fabricación, la integración y los ensayos, planificados
%     para 2026-2027.
%   - AGENTES: rol del practicante (Ingeniero Civil Eléctrico en práctica
%     profesional, elaborador de la documentación SRR y PDR); Juan Carlos
%     Sepúlveda (Ingeniero de Desarrollo Espacial, supervisor directo);
%     Álvaro Valenzuela (PhD en Sistemas Espaciales, supervisión técnica);
%     Lino Acuña (Jefe del Departamento I+D+i, aprobador). Contraparte externa:
%     INVAP (Argentina) como asesor técnico bajo el acuerdo ENAER-DTS-INVAP de
%     marzo de 2025. Equipo I+D+i de 6 personas trabajando en paralelo con el
%     área de diseño mecánico.
%
% [ESTRUCTURA] 4 párrafos:
%   P1  Qué es BIFROST y cuál es su objetivo general (la oración textual de arriba).
%   P2  Cómo funciona técnicamente, en 4-5 oraciones, a nivel de concepto y sin
%       números de diseño (esos van en la sección 2).
%   P3  Plazos, fases cubiertas y alcance/no-alcance. Referencia la Figura 1-1 y
%       la Tabla 1-1.
%   P4  Agentes involucrados y rol concreto del practicante dentro de ellos.
%
% [EJEMPLO DE ARRANQUE DE P3]
%   "El trabajo se desarrolló entre el 2 de enero y el 28 de febrero de 2026 en
%    jornada completa, y cubrió las dos primeras fases del ciclo de vida del
%    proyecto. La Figura 1-1 presenta la carta Gantt maestra elaborada durante la
%    práctica, en la que el área sombreada delimita el período efectivamente
%    ejecutado por el practicante y el resto corresponde a la planificación
%    heredada por el equipo para 2026 y 2027."

% ESCRIBIR AQUÍ

\fig[gantt]{Carta Gantt maestra del proyecto BIFROST. El área
sombreada corresponde al período de práctica profesional (2 de enero al 28 de
febrero de 2026) y los rombos indican los hitos de revisión formal. Elaboración
propia.}{width=\textwidth}{gantt.png}

% [Tabla 1-1] Resume plazos, hitos y entregables. Cubre de una sola vez los
% criterios "plazos" y "alcance" de la rúbrica. Menciónala en el texto de P3.
\begin{table}[H]
\centering
\begin{threeparttable}
\caption{Hitos de revisión y entregables generados durante la práctica.
Elaboración propia.\vspace{1ex}}
\small
\begin{tabularx}{\textwidth}{@{}l l X@{}}
\toprule
\textbf{Fase} & \textbf{Hito y fecha} & \textbf{Entregables generados} \\
\midrule
Pre-Phase A & \textit{Kick-off} & Concepto de Operación (CONOPS) preliminar y
diagrama de bloques inicial del sistema \\
\addlinespace
Phase A & SRR & CONOPS consolidado, 13 requerimientos de usuario, 21
requerimientos de sistema, matriz de trazabilidad, matriz de riesgos y carta
Gantt maestra \\
\addlinespace
Phase B & PDR & Modelo MBSE en Capella con las cuatro capas ARCADIA, herramienta
de simulación de enlace, matriz de ensayos y estimación ROM de costos \\
\bottomrule
\end{tabularx}
\begin{tablenotes}[flushleft]\footnotesize
\item Nota. El \textit{kick-off} se realizó el 16 de enero de 2026, la SRR el 30
de enero de 2026 y la PDR el 27 de febrero de 2026. Las fases posteriores
(Phase C, con cierre en CDR, y Phase D, con cierre en ORR) quedaron fuera del
alcance de la práctica.
\end{tablenotes}
\end{threeparttable}
\end{table}


% ---------------------------------------------------------------------
\subsection{Objetivo de la empresa y problema a resolver}
% ---------------------------------------------------------------------
% [EXTENSIÓN] ~200 palabras.
%
% [CONTENIDO]  (rúbrica B / Identificación, criterio 3)
%   OJO: aquí hay que distinguir DOS objetivos y dejarlos separados en el texto,
%   porque el corrector busca el de la EMPRESA:
%     (a) Objetivo del PROYECTO -> ya lo declaraste en 1.2.
%     (b) Objetivo/interés de la EMPRESA a lograr al finalizar la práctica:
%         - Dar el primer paso formal en el desarrollo de capacidades espaciales
%           dentro de DTS, en el marco del Plan Estratégico 2025-2029, que
%           contempla formar capital humano y adquirir nuevas capacidades.
%         - Llevar BIFROST hasta un nivel de madurez de diseño PDR que demostrara
%           su viabilidad técnica y arquitectónica ANTES de comprometer recursos
%           en diseño detallado y fabricación.
%         - Habilitar la transferencia tecnológica con INVAP en el marco del
%           acuerdo de cooperación, aprovechando la ventana de colaboración.
%   - PROBLEMA / PREGUNTA (formúlalo como pregunta; la rúbrica pide "plantea un
%     problema o pregunta a ser resuelto"):
%       "¿Cómo estructurar y ejecutar, sin un proceso de ingeniería de sistemas
%        preexistente en la empresa, el desarrollo de una carga útil satelital
%        que alcance la madurez de diseño exigida en las revisiones SRR y PDR?"
%
% [ESTRUCTURA] 3 párrafos:
%   P1  El objetivo del proyecto (una línea de recordatorio) y, separado, el
%       objetivo específico de la empresa, con sus tres componentes.
%   P2  Por qué ese objetivo genera un problema de ingeniería y no solo una tarea
%       administrativa: no existía proceso, ni plantillas, ni herencia técnica.
%   P3  Enunciado formal de la pregunta a resolver y anticipo de cómo se aborda.
%
% [EJEMPLO DE ARRANQUE DE P1]
%   "El objetivo específico que DTS definió para el período de práctica no
%    coincide con el objetivo técnico del proyecto. Mientras BIFROST persigue el
%    diseño de una carga útil, el interés de la empresa consistió en materializar
%    su primera incursión formal en el ámbito espacial y en alcanzar, dentro de
%    los dos meses disponibles, un nivel de madurez de diseño suficiente para
%    justificar la continuidad del proyecto ante su jefatura y ante INVAP."

% ESCRIBIR AQUÍ


% ---------------------------------------------------------------------
\subsection{Antecedentes, causas e importancia de su resolución}
% ---------------------------------------------------------------------
% [EXTENSIÓN] ~220 palabras.
%
% [CONTENIDO]  (rúbrica B / Identificación, criterio 4)
%   ANTECEDENTES Y CAUSAS (dos, bien separadas):
%     1. Baja madurez del desarrollo espacial en DTS y en Chile. Consecuencia
%        directa: el proceso de ingeniería, la documentación y las herramientas
%        de soporte debieron generarse sin precedentes internos de referencia.
%        Esta es la causa raíz del problema.
%     2. El acuerdo de cooperación ENAER-DTS-INVAP de marzo de 2025, que habilitó
%        la transferencia de tecnología espacial y de radares y fijó el contexto
%        y el plazo de esta primera incursión.
%   IMPORTANCIA DE RESOLVERLO (tres niveles, de menor a mayor alcance):
%     - Para DTS: consolida experiencia técnica y deja una metodología replicable
%       para futuros proyectos espaciales; evita comprometer capital en un diseño
%       cuya viabilidad no estaba demostrada.
%     - Para el país: se alinea con los objetivos del Centro Espacial Nacional
%       (CEN) y con el impulso de la FACh hacia capacidades espaciales nacionales.
%     - Para el usuario final: Chile tiene geografía extensa y compleja, con zonas
%       aisladas y baja infraestructura terrestre, y alta exposición a sismos y
%       eventos meteorológicos severos; un canal de comunicaciones redundante e
%       independiente de la infraestructura terrestre tiene valor operacional real.
%
% [ESTRUCTURA] 3 párrafos:
%   P1  Antecedente 1 (madurez) + su consecuencia sobre el trabajo.
%   P2  Antecedente 2 (acuerdo INVAP) + la restricción temporal que impone.
%   P3  Importancia en los tres niveles.
%
% [EJEMPLO DE ARRANQUE DE P1]
%   "El problema planteado surge de la confluencia de dos antecedentes. El primero
%    corresponde a la baja madurez del desarrollo espacial, tanto en DTS como en
%    la industria nacional. La empresa no contaba con proyectos espaciales previos,
%    plantillas de documentación ni criterios internos de revisión de diseño, de
%    modo que el practicante debió generar el proceso de ingeniería junto con el
%    producto de ingeniería."

% ESCRIBIR AQUÍ


% =====================================================================
\section{Generación y Justificación}
% =====================================================================
% [EXTENSIÓN TOTAL DE LA SECCIÓN] ~10 planas (1.900 palabras + 5 figuras + 2 tablas).
% Es el corazón del informe:
% vale 4 de los 12 puntos del informe técnico y alimenta la evaluación global.

% ---------------------------------------------------------------------
\subsection{Marco metodológico: ciclo de vida y revisiones de diseño}
% ---------------------------------------------------------------------
% [EXTENSIÓN] ~250 palabras.
%
% [CONTENIDO]  (rúbrica B / Generación, criterio 1: propone una metodología)
%   - El proyecto adoptó un proceso de ingeniería de sistemas secuencial, tipo
%     cascada, estructurado en revisiones formales de madurez del diseño.
%   - Definir qué es una revisión técnica y por qué existe: cada una es una
%     compuerta que exige un conjunto de entregables y criterios de salida antes
%     de autorizar el paso a la fase siguiente.
%     >>> CITA OBLIGATORIA AQUÍ: (Institute of Electrical and Electronics
%         Engineers [IEEE], 2015) -- IEEE Std 15288.2-2014, que define SRR, PDR y
%         CDR con sus criterios de entrada y salida y los entregables exigidos.
%   - Explicar las tres compuertas relevantes:
%       SRR -> la línea base de requerimientos está completa y es verificable.
%       PDR -> la arquitectura propuesta es técnicamente viable y cumple los
%              requerimientos; se autoriza el diseño detallado.
%       CDR -> el diseño está listo para fabricación (fuera del alcance).
%   - Justificar la elección de cascada frente a un enfoque iterativo: el hardware
%     espacial no admite iteración barata después de la fabricación, el costo del
%     cambio crece con la fase, y el proyecto tenía una contraparte externa que
%     revisaba en hitos discretos.
%
% [ESTRUCTURA] 3 párrafos:
%   P1  Qué metodología se adoptó y por qué el ciclo de vida secuencial es el
%       adecuado para hardware espacial (justificación, no solo descripción).
%   P2  Qué es una revisión de madurez y qué exige cada una (con la cita IEEE).
%   P3  Qué fases cubrió la práctica y hasta dónde llegó (enlazar con Figura 1-1).
%
% [EJEMPLO DE ARRANQUE DE P2]
%   "Cada revisión funciona como una compuerta de madurez y no como una simple
%    presentación de avance. El estándar IEEE Std 15288.2-2014 establece, para
%    cada una de ellas, los criterios de entrada, los criterios de salida y el
%    conjunto de entregables que el equipo debe presentar (IEEE, 2015). Bajo ese
%    marco, la Revisión de Requerimientos del Sistema exige una línea base de
%    requerimientos completa, no ambigua y verificable, mientras que la Revisión
%    de Diseño Preliminar exige demostrar que la arquitectura propuesta satisface
%    esos requerimientos antes de comprometer recursos en el diseño detallado."

% ESCRIBIR AQUÍ


% ---------------------------------------------------------------------
\subsection{Concepto de Operación}
% ---------------------------------------------------------------------
% [EXTENSIÓN] ~300 palabras + Figura 2-1.
%
% [CONTENIDO]
%   - Qué es un CONOPS y por qué fue el PRIMER entregable: describe el sistema
%     desde la perspectiva del usuario y fija la línea base operativa de la que
%     se derivan todos los requerimientos.
%   - Desarrolla las seis preguntas del CONOPS, pero EN PROSA, no como lista:
%       QUÉ:     carga útil Store & Forward en plataforma CubeSat 12U, repetidor
%                no instantáneo de mensajes digitales provenientes de estaciones
%                remotas autónomas (DCP).
%       POR QUÉ: conectividad en zonas de geografía compleja, resiliencia ante
%                emergencias, plataforma experimental con proyección a IoT satelital.
%       QUIÉNES: DTS (desarrollador integral), INVAP (contraparte técnica);
%                usuarios potenciales: universidades y centros de investigación,
%                centros de monitoreo meteorológico, sismológico y vulcanológico,
%                organismos de gestión de emergencias, empresas de IoT satelital.
%       DÓNDE:   tres segmentos -- espacial (LEO SSO, 650 km, 96 grados),
%                terrestre (Chile continental, zonas extremas, insular y
%                antártico) y espectral (uplink UHF 401-403 MHz, downlink UHF
%                ~466 MHz).
%       CUÁNDO:  vida útil de 5 años, 4 a 5 pasadas diarias sobre territorio
%                nacional, ventanas de contacto del orden de 10 minutos.
%       CÓMO:    las DCP transmiten de forma autónoma; la carga útil escucha
%                dentro de las ventanas de misión, compensa Doppler, demodula y
%                valida; los mensajes válidos se temporalizan y almacenan en
%                memoria no volátil; la descarga ocurre bajo comando del OBC.
%   - Menciona el estándar de referencia para la interoperabilidad de las
%     estaciones remotas: el sistema Argos de recolección de datos ambientales
%     por satélite.  >>> CITA: (CLS, 2016).
%
% [ESTRUCTURA] 4 párrafos:
%   P1  Qué es un CONOPS, por qué fue el primer entregable y qué rol cumple como
%       línea base (esto justifica la decisión metodológica).
%   P2  Qué es el sistema y por qué se desarrolla.
%   P3  Dónde y cuándo opera: los tres segmentos y el marco orbital.
%   P4  Cómo opera: la cadena de operación en cinco pasos, redactada en prosa.
%       Referencia la Figura 2-1.
%
% [EJEMPLO DE ARRANQUE DE P4]
%   "La operación nominal del sistema encadena cinco etapas. Las estaciones
%    remotas transmiten sus mensajes de forma autónoma en la banda de 401 a 403
%    MHz, sin coordinación previa con el satélite. La carga útil escucha
%    activamente dentro de las ventanas definidas por la misión y compensa el
%    corrimiento en frecuencia inducido por el efecto Doppler, que alcanza
%    $\pm$10 kHz debido a la velocidad relativa del satélite, cercana a 7,5 km/s."

% ESCRIBIR AQUÍ

\fig[entidades]{Entidades operacionales del sistema BIFROST y sus
relaciones, modeladas en Eclipse Capella durante la fase de análisis operacional.
Elaboración propia.}{width=0.95\textwidth}{oeb_entidades_operacionales.jpg}

% [FIGURA OPCIONAL, si te sobra espacio] Escenario operacional de recolección:
% \fig[escenario]{Escenario operacional de recolección de datos desde
% estaciones remotas hacia la carga útil. Elaboración propia.}%
% {width=0.95\textwidth}{oab_recoleccion_datos.jpg}


% ---------------------------------------------------------------------
\subsection{Definición y trazabilidad de requerimientos}
% ---------------------------------------------------------------------
% [EXTENSIÓN] ~350 palabras + Tabla 2-1 + Tabla 2-2.
%
% [CONTENIDO]  (rúbrica B / Generación, criterios 1 y 4)
%   - Proceso: del CONOPS se derivaron 13 requerimientos de usuario (UR) y de
%     estos, 21 requerimientos de sistema (SR), numerados hasta SR-23 por
%     depuraciones sucesivas de la lista. Explica la lógica de la derivación:
%     el UR expresa una necesidad, el SR la traduce a una especificación técnica
%     verificable.
%   - Matriz de trazabilidad bidireccional: cada UR tiene al menos un SR asociado
%     y cada SR responde a al menos un UR. Explica para qué sirve: permite
%     detectar requerimientos huérfanos y evaluar el impacto de un cambio.
%   - Atributos asignados a cada requerimiento: prioridad, riesgo, dificultad,
%     tipo (funcional / no funcional / interfaz / proceso) y estado.
%   - Restricciones que acotaron el diseño (Tabla 2-2): potencia, masa, centro de
%     gravedad, volumen, radiación acumulada, desgasificación y rango térmico.
%     Explica de dónde vienen: la plataforma CubeSat 12U, el lanzador y el entorno
%     orbital, no una decisión del equipo.
%   - Evaluación de riesgo por requerimiento: impacto y probabilidad en escala de
%     tres niveles, producto de ambos como valor de riesgo. Menciona el resultado
%     agregado y remite al Anexo C.
%
% [ESTRUCTURA] 4 párrafos:
%   P1  Cómo se derivaron los UR desde el CONOPS y qué caracteriza a un buen
%       requerimiento (verificable, no ambiguo, único). Referencia la Tabla 2-1.
%   P2  Cómo se derivaron los SR y qué asegura la matriz de trazabilidad.
%   P3  Las restricciones de plataforma y entorno. Referencia la Tabla 2-2 y
%       explica su efecto: acotaron el espacio de soluciones antes de diseñar.
%   P4  La evaluación de riesgo, el criterio usado y qué se hizo con los riesgos
%       altos (mitigaciones tempranas). Remite al Anexo C.
%
% [EJEMPLO DE ARRANQUE DE P2]
%   "Cada requerimiento de usuario se tradujo luego a uno o más requerimientos de
%    sistema, redactados en términos medibles y asociados a un método de
%    verificación. El requerimiento UR-01, que exige un servicio Store & Forward
%    con almacenamiento durante la órbita, derivó en los requerimientos SR-04 y
%    SR-06, que fijan una capacidad mínima de 1.000.000 de mensajes retenidos
%    durante hasta 15 órbitas y una descarga de al menos 10.000 mensajes por
%    ventana de contacto. Esta traducción es la que hace verificable una
%    necesidad que, en su forma original, no lo era."

% ESCRIBIR AQUÍ

\begin{table}[H]
\centering
\begin{threeparttable}
\caption{Requerimientos de usuario clasificados con nivel de riesgo
alto y sus medidas de mitigación tempranas. Elaboración propia a partir de la
matriz de trazabilidad del proyecto.\vspace{1ex}}
\small
\begin{tabularx}{\textwidth}{@{}l X c X@{}}
\toprule
\textbf{ID} & \textbf{Requerimiento de usuario} & \textbf{Riesgo} & \textbf{Mitigación definida} \\
\midrule
UR-01 & Proveer un servicio \textit{Store \& Forward} satelital, almacenando datos
durante la órbita para descargarlos en una estación terrestre & Alto &
Diseño de memoria no volátil con margen $\geq$ 50\,\%, simulación de carga orbital
y pruebas de estrés en tierra \\
\addlinespace
UR-06 & Ser compatible con una plataforma CubeSat 12U en órbita LEO sol-síncrona
a 650 km de altitud y 96$^{\circ}$ de inclinación & Alto &
Análisis de interfaces mecánicas y eléctricas, y revisión conjunta con el
proveedor de la plataforma \\
\addlinespace
UR-08 & Permitir el uso de componentes COTS con funcionamiento en órbita probado,
tolerancia a la radiación verificada o mitigaciones de diseño & Alto &
Selección de componentes con validación espacial previa y mitigación por
redundancia y \textit{watchdogs} \\
\addlinespace
UR-11 & Garantizar una vida útil operativa mínima de 5 años en el entorno de
radiación LEO definido & Alto &
Pruebas de estrés, análisis térmico y pruebas de vida aceleradas \\
\bottomrule
\end{tabularx}
\begin{tablenotes}[flushleft]\footnotesize
\item Nota. La lista completa de los 13 requerimientos de usuario se presenta en
el Anexo A y la de los 21 requerimientos de sistema en el Anexo B.
\end{tablenotes}
\end{threeparttable}
\end{table}

\begin{table}[H]
\centering
\begin{threeparttable}
\caption{Restricciones impuestas por la plataforma CubeSat 12U y por el
entorno orbital sobre el diseño de la carga útil. Elaboración propia a partir de
la documentación del proyecto.\vspace{1ex}}
\small
\begin{tabularx}{\textwidth}{@{}l l X@{}}
\toprule
\textbf{Parámetro} & \textbf{Límite} & \textbf{Origen de la restricción} \\
\midrule
Potencia            & 10 W en operación nominal & Sistema de potencia de la plataforma \\
Masa                & 3,0 kg (máximo)           & Presupuesto de masa del CubeSat 12U \\
Centro de gravedad  & $\pm$2 mm respecto del centro geométrico & Estabilidad del control de actitud \\
Volumen             & 20 $\times$ 20 $\times$ 10 cm & Envolvente asignada a la carga útil \\
Dosis acumulada     & 10 a 20 krad en 5 años    & Ambiente radiativo en órbita LEO \\
Desgasificación     & TML $<$ 1,0\,\% y CVCM $<$ 0,1\,\% & Protección de sensores ópticos y paneles solares \\
Rango térmico       & $-20$ a $+60\ ^{\circ}$C en operación & Ciclado térmico orbital \\
Efecto Doppler      & $\pm$10 kHz                & Velocidad relativa del satélite ($\sim$7,5 km/s) \\
\bottomrule
\end{tabularx}
\end{threeparttable}
\end{table}


% ---------------------------------------------------------------------
\subsection{Modelamiento MBSE con ARCADIA y Eclipse Capella}
% ---------------------------------------------------------------------
% [EXTENSIÓN] ~450 palabras en total (intro 130, 2.4.1 100, 2.4.2 90, 2.4.3 130),
% más tres figuras.
%
% [CONTENIDO GENERAL]  (rúbrica B / Generación, criterios 2 y 3: describe los
%  softwares necesarios y justifica su elección)
%   - Punto de partida: el proyecto arrancó con documentación y diagramas de
%     bloques sueltos. Al crecer el número de interfaces (plataforma, estaciones
%     remotas, estaciones terrenas, centro de control de misión), mantener la
%     consistencia entre documentos se volvió inviable.
%   - DECISIÓN: migrar el proyecto completo a MBSE (Model-Based Systems
%     Engineering) con la metodología ARCADIA, modelada en Eclipse Capella.
%     >>> Esta migración es TUYA y ocurrió a mitad de la práctica (semana 6 de la
%         bitácora). Dilo explícitamente: es el mejor argumento para el criterio
%         de "aplicación creativa y/o novedosa dentro de la empresa" (rúbrica F).
%   - JUSTIFICACIÓN de la decisión (no basta con describirla):
%       * el modelo pasa a ser la única fuente de verdad, en vez de N documentos;
%       * la trazabilidad entre función, componente lógico y componente físico es
%         automática, no manual;
%       * Capella es de código abierto, sin costo de licencia, lo que importaba
%         en un proyecto exploratorio sin presupuesto asignado;
%       * ARCADIA está ampliamente adoptada en desarrollo de sistemas complejos
%         y aeroespaciales.  >>> CITA: (Roques, 2018).
%   - Las cuatro fases de ARCADIA: análisis operacional (OA), análisis del sistema
%     (SA), arquitectura lógica (LA) y arquitectura física (PA).
%
% [ESTRUCTURA DE LA INTRODUCCIÓN] 2 párrafos antes del primer \subsubsection:
%   P1  El problema de consistencia que apareció y por qué la documentación
%       tradicional no lo resolvía.
%   P2  Qué es MBSE/ARCADIA, por qué se eligió Capella y cuáles son sus cuatro
%       fases (con la cita).
%
% [EJEMPLO DE ARRANQUE DE P1]
%   "Hacia la mitad del período de práctica, el proyecto acumulaba un CONOPS, dos
%    matrices de requerimientos y varios diagramas de bloques independientes entre
%    sí. Cada modificación de un requerimiento obligaba a revisar manualmente los
%    demás documentos para mantener la coherencia, un procedimiento que no escala
%    en un sistema con cuatro interfaces externas. Frente a esa limitación, el
%    practicante propuso migrar el proyecto completo a un enfoque de ingeniería de
%    sistemas basada en modelos."

% ESCRIBIR AQUÍ (introducción de la subsección)

\subsubsection{Análisis operacional y análisis del sistema}
% [EXTENSIÓN] ~100 palabras + Figura 2-2.
%
% [CONTENIDO]
%   - OA: define qué deben lograr los usuarios y el entorno de la misión, sin
%     hablar todavía del sistema. Entidades operacionales: estaciones remotas,
%     centro de control de misión, estaciones terrenas, plataforma satelital.
%   - SA: trata a BIFROST como una caja negra. Define sus funciones de frontera y
%     los intercambios de datos con los actores externos. Aquí aparecen las
%     capacidades del sistema y los escenarios (recolección de datos Argos-2 y
%     Argos-3, actualización de firmware).
%   - Explica la diferencia entre ambos niveles en una oración: el OA describe el
%     problema; el SA describe qué debe hacer el sistema, todavía sin decir cómo.
%
% [ESTRUCTURA] 2 párrafos, uno por nivel. Referencia la Figura 2-2 en el segundo.

% ESCRIBIR AQUÍ

\fig[sab]{Arquitectura del sistema BIFROST tratado como caja negra,
con sus funciones de frontera e intercambios de datos con los actores externos.
Elaboración propia.}{width=0.95\textwidth}{sab_sistema_bifrost.jpg}

\subsubsection{Arquitectura lógica}
% [EXTENSIÓN] ~90 palabras + Figura 2-3.
%
% [CONTENIDO]
%   - La LA abre la caja negra y agrupa las funciones en componentes lógicos,
%     independientes de la tecnología con que se implementen.
%   - Los seis componentes lógicos definidos: RF Frontend, Digital Processing,
%     Data Handling, Power Management, Interface Handling y Data Storage.
%   - Explica qué hace cada uno en media línea y cómo se encadenan en la cadena
%     funcional de recolección de datos (RF Frontend -> Digital Processing ->
%     Data Handling -> Data Storage) y en la de transmisión (sentido inverso).
%   - Por qué importa este nivel intermedio: permite decidir la arquitectura
%     funcional antes de comprometerse con componentes concretos, y hace que un
%     cambio de proveedor no obligue a rehacer el diseño funcional.
%
% [ESTRUCTURA] 2 párrafos:
%   P1  Qué es la arquitectura lógica y cuáles son los seis componentes.
%   P2  Cómo se encadenan en las dos cadenas funcionales y qué ventaja aporta el
%       desacople respecto de la tecnología. Referencia la Figura 2-3.

% ESCRIBIR AQUÍ

\fig[lcbd]{Descomposición lógica del sistema en sus seis componentes
principales, independiente de la implementación tecnológica. Elaboración propia.}%
{width=0.9\textwidth}{lcbd_arquitectura_logica.jpg}

\subsubsection{Arquitectura física}
% [EXTENSIÓN] ~130 palabras + Figura 2-4.
%
% [CONTENIDO]  (este es el resultado técnico más fuerte de la práctica)
%   - La PA asigna las funciones y los componentes lógicos a hardware y software
%     concretos.
%   - COMPONENTE CLAVE PRESELECCIONADO: la plataforma SDR TOTEM de Alén Space,
%     que integra el transceptor AD9364 y un SoC Zynq-7020 (FPGA + procesador
%     Cortex-A9).  >>> CITA: (Alén Space, s.f.)
%     Reparto de funciones:
%       * AD9364      -> filtrado de LO/banda base y conversión ADC/DAC;
%       * FPGA        -> análisis de frecuencia, seguimiento Doppler, filtrado
%                        adaptativo, demodulación y modulación;
%       * Cortex-A9   -> verificación de integridad, identificación de metadatos
%                        de la DCP, gestión de memoria e interfaz de bus CAN.
%   - Bloques desarrollados a medida por DTS: el frontend RF (switch T/R,
%     amplificador de potencia y LNA), el subsistema de antenas y la red de
%     distribución de potencia con detección de latch-up, monitoreo de consumo y
%     reinicio por watchdog.
%   - JUSTIFICA la decisión COTS + SDR:
%       * un SDR con herencia de vuelo reduce el riesgo asociado a UR-08;
%       * la demodulación en FPGA permite adaptar el esquema de modulación por
%         software sin rediseñar hardware, lo que responde al requerimiento de
%         escalabilidad (UR-10) y a la actualización en vuelo (SR-17);
%       * concentrar el desarrollo propio en el frontend RF y la gestión de
%         potencia acota el esfuerzo de ingeniería al tiempo disponible.
%   - Mitigaciones de radiación en la arquitectura: redundancia de procesador y
%     bancos de memoria, algoritmos CRC/EDAC y watchdog timers externos.
%
% [ESTRUCTURA] 3 párrafos:
%   P1  Qué asigna la arquitectura física y cuál es el componente central elegido.
%   P2  Reparto de funciones entre AD9364, FPGA y Cortex-A9, y qué desarrolla DTS
%       a medida. Referencia la Figura 2-4.
%   P3  Justificación de la decisión y mitigaciones de radiación incorporadas.
%
% [EJEMPLO DE ARRANQUE DE P3]
%   "La selección de una plataforma SDR comercial con herencia de vuelo responde
%    directamente a dos requerimientos en tensión. Por una parte, UR-08 admite
%    componentes COTS solo si acreditan funcionamiento en órbita, tolerancia a la
%    radiación verificada o mitigaciones de diseño equivalentes. Por otra, UR-10
%    exige una arquitectura modular y escalable. Un receptor implementado en
%    lógica programable satisface ambas condiciones: hereda la calificación del
%    módulo comercial y permite modificar el esquema de demodulación por software,
%    sin rediseñar el hardware ni repetir la campaña de ensayos ambientales."

% ESCRIBIR AQUÍ

\fig[pab]{Arquitectura física de la carga útil BIFROST. La cadena de
recolección de datos se destaca en rojo y la cadena de transmisión en azul.
Elaboración propia.}{width=\textwidth}{pab_arquitectura_fisica.jpg}

% [ALTERNATIVA] Si el diagrama queda ilegible en vertical, usa una página
% apaisada. El paquete lscape ya está cargado en style.cls. Comenta el \fig de
% arriba y descomenta este bloque:
% \begin{landscape}
% \fig[pab]{Arquitectura física de la carga útil BIFROST. La cadena de
% recolección de datos se destaca en rojo y la cadena de transmisión en azul.
% Elaboración propia.}{width=\linewidth}{pab_arquitectura_fisica.jpg}
% \end{landscape}


% ---------------------------------------------------------------------
\subsection{Herramienta de simulación orbital y presupuesto de enlace}
% ---------------------------------------------------------------------
% [EXTENSIÓN] ~280 palabras + Tabla 2-3 + Figura 2-5.
%
% [CONTENIDO]  (rúbrica B / Generación, criterio 2: describe los análisis,
%               mediciones y softwares necesarios para resolver el problema)
%   - PROBLEMA A RESPONDER: ninguna herramienta disponible en el departamento
%     permitía evaluar simultáneamente la geometría orbital y el presupuesto de
%     enlace considerando los patrones de radiación reales de ambas antenas.
%   - SOLUCIÓN: el practicante desarrolló una herramienta propia en Python que
%     propaga una órbita sol-síncrona en tres dimensiones y calcula el presupuesto
%     de enlace en cada instante.
%     >>> SEGUNDO ARGUMENTO FUERTE PARA "aplicación creativa" (rúbrica F).
%   - Arquitectura del código (la rúbrica pide "los softwares"):
%       * módulo de mecánica orbital: propagación de la órbita y geometría
%         satélite-estación (distancia de enlace y ángulo de elevación);
%       * módulo de presupuesto de enlace: pérdidas de espacio libre, pérdidas
%         adicionales y relación Eb/N0 resultante;
%       * módulos de patrón de radiación: monopolo con efecto de tierra para la
%         estación terrestre y turnstile para el satélite;
%       * archivo central de configuración y script de visualización interactiva.
%   - PARÁMETROS DE SIMULACIÓN -> van en la Tabla 2-3, no en el texto.
%   - Explica por qué modelar el patrón de radiación cambia el resultado: un
%     cálculo de enlace que asume ganancia constante sobrevalora el margen en los
%     extremos del patrón y no revela los nulos.
%   - Justifica la elección de Python y del desarrollo propio frente a software
%     comercial: costo nulo, control total sobre el modelo de antena, código
%     reutilizable por la empresa en proyectos posteriores.
%
% [ESTRUCTURA] 4 párrafos:
%   P1  Qué análisis había que hacer y por qué las herramientas disponibles no
%       bastaban.
%   P2  Qué se construyó y cómo está organizado el código.
%   P3  Qué parámetros y modelos de antena se usaron. Referencia la Tabla 2-3.
%   P4  Justificación de la decisión y anticipo de que los resultados se analizan
%       en la sección 3.
%
% [EJEMPLO DE ARRANQUE DE P1]
%   "La viabilidad de BIFROST depende de una condición que ninguna decisión de
%    arquitectura puede compensar: que el enlace cierre. Verificarlo exigía
%    evaluar de forma conjunta la geometría orbital, que determina la distancia y
%    el ángulo de elevación en cada instante, y el balance de potencias, que
%    depende de la ganancia que cada antena presenta precisamente en esa
%    dirección. El departamento no disponía de una herramienta que acoplara ambos
%    cálculos, por lo que el practicante desarrolló un entorno de simulación
%    propio en Python."

% ESCRIBIR AQUÍ

\begin{table}[H]
\centering
\begin{threeparttable}
\caption{Parámetros empleados en la simulación del presupuesto de
enlace. Elaboración propia.\vspace{1ex}}
\small
\begin{tabularx}{\textwidth}{@{}X l@{}}
\toprule
\textbf{Parámetro} & \textbf{Valor} \\
\midrule
Altitud orbital nominal                       & 650 km \\
Tipo de órbita e inclinación                  & Sol-sincrónica (SSO), 96$^{\circ}$ \\
Frecuencia de enlace de subida                & 401--403 MHz \\
Frecuencia de enlace de bajada                & $\sim$466 MHz \\
Potencia de transmisión                       & 1 W (0 dBW) \\
Antena de la estación terrestre               & Monopolo con efecto de plano de tierra \\
Antena del satélite                           & \textit{Turnstile} \\
Pérdidas por desacople de polarización        & 3 dB \\
Pérdidas por desadaptación de impedancia      & 2 dB \\
Umbral de Eb/N0 para cierre de enlace         & 14 dB \\
Tasa de datos                                 & 400 bps \\
\bottomrule
\end{tabularx}
\begin{tablenotes}[flushleft]\footnotesize
\item Nota. Las ganancias de antena corresponden a los valores típicos de cada
configuración, modeladas a partir de su patrón de radiación completo y no como
una ganancia constante.
\end{tablenotes}
\end{threeparttable}
\end{table}

\fig[uplinksim]{Interfaz de la herramienta de simulación desarrollada
en Python. El panel superior presenta la relación Eb/N0 del enlace de subida a lo
largo de dos días de propagación orbital, con el umbral de cierre de enlace
indicado en línea segmentada. Elaboración propia.}{width=\textwidth}{uplink.png}


% ---------------------------------------------------------------------
\subsection{Justificación de las decisiones metodológicas y de recursos}
% ---------------------------------------------------------------------
% [EXTENSIÓN] ~280 palabras.
%
% [CONTENIDO]  (rúbrica B / Generación, criterio 3. Es un criterio propio: no
%  basta con haber justificado cosas sueltas antes, conviene un cierre explícito.)
%   Recoge las cuatro decisiones mayores y justifica cada una contra DOS ejes:
%   (a) las características del problema y (b) las características de DTS.
%
%   D1. Ciclo de vida secuencial con revisiones formales
%       Problema: hardware espacial, costo del cambio creciente, no hay iteración
%                 barata después del lanzamiento.
%       Empresa:  DTS ya opera bajo ISO 9001:2015 y certificaciones aeronáuticas,
%                 de modo que la cultura de revisiones documentadas ya existía.
%
%   D2. MBSE con ARCADIA sobre Eclipse Capella
%       Problema: múltiples interfaces externas, consistencia inmanejable a mano.
%       Empresa:  herramienta de código abierto, sin costo de licencia en un
%                 proyecto exploratorio; deja un modelo reutilizable y un método
%                 replicable para los siguientes proyectos espaciales de DTS.
%
%   D3. Componentes COTS con herencia de vuelo sobre una plataforma SDR
%       Problema: dosis de radiación, plazos y presupuesto acotado.
%       Empresa:  DTS no cuenta con capacidad de fabricación de electrónica de
%                 grado espacial; concentrar el desarrollo propio en el frontend
%                 RF aprovecha su experiencia real en electrónica y aviónica.
%
%   D4. Herramienta de simulación propia en Python
%       Problema: necesidad de acoplar geometría orbital y patrón de antena.
%       Empresa:  sin licencias comerciales disponibles; el código queda como
%                 activo interno reutilizable.
%
%   - CIERRA con las limitaciones metodológicas asumidas conscientemente: el
%     modelo de enlace no incorpora pérdidas atmosféricas ni interferencia
%     co-canal, y la selección de componentes es preliminar y sujeta a
%     confirmación en CDR.
%
% [ESTRUCTURA] 5 párrafos: uno por decisión (D1 a D4) y uno de cierre sobre
% limitaciones asumidas. Cada párrafo con el mismo patrón: decisión -> por qué
% el problema la exige -> por qué el contexto de DTS la hace la opción correcta.
%
% [EJEMPLO DE ARRANQUE DE D2]
%   "La migración a un enfoque basado en modelos respondió tanto a la naturaleza
%    del sistema como a las condiciones de la empresa. Un sistema con cuatro
%    interfaces externas y dos cadenas funcionales acopladas exige un mecanismo
%    de trazabilidad que la documentación en texto no ofrece. Al mismo tiempo,
%    Eclipse Capella es una herramienta de código abierto, lo que resultaba
%    determinante en un proyecto exploratorio sin presupuesto de licencias
%    asignado, y deja instalada en DTS una metodología reutilizable en los
%    proyectos espaciales posteriores."

% ESCRIBIR AQUÍ


% =====================================================================
\section{Evaluación y Conclusiones}
% =====================================================================
% [EXTENSIÓN TOTAL DE LA SECCIÓN] ~5,5 planas (1.300 palabras + 2 figuras + 2 tablas).
% ATENCIÓN: esta sección vale otros 4 puntos y es donde más informes pierden
% puntaje, porque describen actividades en vez de resultados medibles.

% ---------------------------------------------------------------------
\subsection{Resultados obtenidos}
% ---------------------------------------------------------------------
% [EXTENSIÓN] ~350 palabras + Tabla 3-1 + Figuras 3-1 y 3-2 + Tabla 3-2.
%
% [CONTENIDO]  (rúbrica B / Evaluación, criterio 1: "resultados apoyados en
%               indicadores concretos o mediciones reales o estimadas")
%   >>> REGLA: cada afirmación de esta subsección debe llevar un número.
%
%   R1. Resultado de gestión de la ingeniería (Tabla 3-1):
%       - 13 requerimientos de usuario y 21 de sistema definidos y trazados;
%       - 35 relaciones de trazabilidad; cobertura del 100 % en AMBOS sentidos
%         (ningún UR sin SR asociado y ningún SR huérfano);
%       - 12 requerimientos clasificados en nivel de riesgo alto (4 UR y 8 SR),
%         todos con mitigación definida;
%       - 19 ensayos de verificación especificados (7 funcionales, 3 térmicos,
%         5 mecánicos y 4 de compatibilidad electromagnética);
%       - modelo MBSE con las cuatro capas ARCADIA completas.
%
%   R2. Resultado de viabilidad del enlace (Figuras 2-5, 3-1 y 3-2):
%       - el enlace cierra en 4 a 5 pasadas diarias sobre el territorio nacional;
%       - ventanas de contacto útiles del orden de 10 minutos por pasada;
%       - distancia media de enlace en condición de contacto: ~1.500 km;
%       - Eb/N0 por sobre el umbral de 14 dB durante toda la ventana, con
%         máximos en torno a 30-35 dB;
%       - HALLAZGO NO TRIVIAL, DESTÁCALO: la simulación revela un nulo profundo
%         en el cénit, cuando la elevación se acerca a 90 grados y la distancia
%         es mínima (655 km). Es consecuencia directa del patrón del monopolo y
%         no aparecería en un cálculo con ganancia constante. Es exactamente el
%         tipo de hallazgo que justifica haber construido la herramienta.
%
%   R3. Resultado programático (Tabla 3-2):
%       - estimación ROM de 133.000 USD para entregar un modelo de vuelo integrado;
%       - cronograma con las fases hasta la revisión AIT (Figura 1-1).
%
% [ESTRUCTURA] 4 párrafos:
%   P1  Resultados de la ingeniería de requerimientos, con los números. Referencia
%       la Tabla 3-1.
%   P2  Resultados del presupuesto de enlace: cuántas pasadas, qué duración, qué
%       márgenes. Referencia la Figura 3-1.
%   P3  El hallazgo del nulo cenital y qué implica para la operación (la ventana
%       útil no es continua; conviene planificar la descarga fuera del cénit).
%       Referencia la Figura 3-2.
%   P4  Resultado programático: cronograma y estimación de costos. Referencia la
%       Tabla 3-2.
%
% [EJEMPLO DE ARRANQUE DE P3]
%   "La simulación reveló además un comportamiento que un cálculo de enlace con
%    ganancia constante no habría detectado. En las pasadas de alta elevación, la
%    relación Eb/N0 cae abruptamente por debajo del umbral en el instante de
%    máxima aproximación, pese a que en ese punto la distancia de enlace alcanza
%    su valor mínimo de 655 km. La causa es el nulo cenital del patrón de
%    radiación del monopolo de la estación terrestre, que anula la ganancia
%    justamente en la dirección de menor pérdida de espacio libre."

% ESCRIBIR AQUÍ

\begin{table}[H]
\centering
\begin{threeparttable}
\caption{Indicadores de los entregables generados durante el período de
práctica. Elaboración propia.\vspace{1ex}}
\small
\begin{tabularx}{\textwidth}{@{}X l l@{}}
\toprule
\textbf{Indicador} & \textbf{Resultado} & \textbf{Estado} \\
\midrule
Requerimientos de usuario definidos                       & 13    & Línea base cerrada en SRR \\
Requerimientos de sistema derivados                       & 21    & Línea base cerrada en SRR \\
Relaciones de trazabilidad UR $\leftrightarrow$ SR        & 35    & Matriz bidireccional \\
Cobertura de trazabilidad en ambos sentidos               & 100\,\% & Verificada \\
Requerimientos evaluados con riesgo alto                  & 12    & Mitigaciones propuestas \\
Ensayos de verificación especificados                     & 19    & Matriz preliminar \\
Capas del modelo ARCADIA completadas                      & 4 de 4 & Modelo entregado \\
Revisiones formales cerradas                              & 2 (SRR y PDR) & Aprobadas con INVAP \\
\bottomrule
\end{tabularx}
\begin{tablenotes}[flushleft]\footnotesize
\item Nota. Los 19 ensayos se distribuyen en 7 funcionales, 3 térmicos, 5
mecánicos y 4 de compatibilidad electromagnética. El detalle se presenta en el
Anexo D.
\end{tablenotes}
\end{threeparttable}
\end{table}

\fig[downlinksim]{Relación Eb/N0 del enlace de bajada a lo largo de
dos días de propagación orbital. Cada pico corresponde a una pasada sobre la
estación terrena y la línea segmentada indica el umbral de cierre de enlace.
Elaboración propia.}{width=\textwidth}{downlink.png}

\fig[ventana]{Detalle de una ventana de contacto individual. La caída
abrupta en el centro de la ventana corresponde al nulo cenital del patrón de
radiación de la antena terrestre. Elaboración propia.}%
{width=\textwidth}{zoom_uplink.png}

\begin{table}[H]
\centering
\begin{threeparttable}
\caption{Estimación de costos de orden de magnitud (ROM) para la
entrega de un modelo de vuelo integrado a la plataforma satelital. Elaboración
propia.\vspace{1ex}}
\small
\begin{tabularx}{\textwidth}{@{}l X r@{}}
\toprule
\textbf{Ítem} & \textbf{Descripción} & \textbf{Costo estimado (USD)} \\
\midrule
Recursos humanos          & Dos ingenieros durante 12 meses                & 48.000 \\
Componentes y materiales  & Electrónica, mecánica y materiales             & 50.000 \\
AIT                       & Integración, pruebas funcionales y ambientales & 15.000 \\
Logística                 & Transporte, envíos, importaciones y soporte    & 20.000 \\
\midrule
\textbf{Total}            &                                                & \textbf{133.000} \\
\bottomrule
\end{tabularx}
\begin{tablenotes}[flushleft]\footnotesize
\item Nota. La estimación asume el uso de componentes y materiales comerciales
disponibles en el mercado y la ejecución de las actividades de integración,
ensamblado y pruebas en instalaciones nacionales.
\end{tablenotes}
\end{threeparttable}
\end{table}


% ---------------------------------------------------------------------
\subsection{Análisis, limitaciones e implicancias futuras}
% ---------------------------------------------------------------------
% [EXTENSIÓN] ~300 palabras.
%
% [CONTENIDO]  (rúbrica B / Evaluación, criterio 2: limitaciones y relevancia de
%  los análisis, implicancias futuras de las decisiones y pasos a seguir)
%   ANÁLISIS -- qué significan los resultados:
%     - El cierre de enlace con margen sobre el umbral en 4 a 5 pasadas diarias
%       confirma que el concepto operativo es viable con la potencia y las antenas
%       supuestas, y que la restricción real no es el enlace sino el tiempo
%       acumulado de contacto.
%     - La arquitectura basada en SDR resuelve simultáneamente los requerimientos
%       de tolerancia a radiación, modularidad y actualización en vuelo.
%
%   LIMITACIONES -- sé explícito, la rúbrica lo pide y reconocerlas suma:
%     - El modelo de enlace no incorpora pérdidas atmosféricas, ruido ambiental ni
%       interferencia co-canal con otros usuarios de la banda Argos.
%     - Se asume una única estación terrestre y una única DCP; el escenario real
%       con múltiples estaciones transmitiendo simultáneamente introduce colisiones
%       que el modelo actual no representa.
%     - La preselección de componentes es preliminar: no se verificó la dosis
%       acumulada real ni el presupuesto de potencia detallado por modo operativo.
%     - Los 19 ensayos están especificados pero ninguno ejecutado; los valores de
%       sensibilidad (-115 dBm) y consumo son estimaciones de diseño.
%     - La estimación de costos es de orden de magnitud, con incertidumbre alta.
%
%   IMPLICANCIAS FUTURAS DE LAS DECISIONES TOMADAS:
%     - Adoptar Capella obliga a mantener el modelo actualizado; si el equipo
%       vuelve a documentar en texto, se pierde la trazabilidad ganada.
%     - Elegir una plataforma SDR comercial crea dependencia de un proveedor
%       único y de su cadena de suministro internacional.
%     - El nulo cenital detectado condiciona la planificación de misión y obliga
%       a revisar el patrón de la antena terrestre en la fase siguiente.
%
%   PASOS A SEGUIR (concretos, no genéricos):
%     - Incorporar pérdidas atmosféricas y un escenario multi-DCP al simulador;
%     - cerrar el presupuesto de potencia por modo operativo y el análisis térmico;
%     - confirmar la selección de componentes y levantar la BOM definitiva;
%     - ejecutar la campaña de ensayos funcionales y ambientales en el CEN;
%     - avanzar a la Revisión Crítica de Diseño (CDR) prevista para abril de 2026.
%
% [ESTRUCTURA] 4 párrafos: análisis, limitaciones, implicancias, pasos a seguir.
% Un párrafo por bloque, sin viñetas.

% ESCRIBIR AQUÍ


% ---------------------------------------------------------------------
\subsection{Cumplimiento del objetivo de la empresa}
% ---------------------------------------------------------------------
% [EXTENSIÓN] ~220 palabras.
%
% [CONTENIDO]  (rúbrica B / Evaluación, criterio 3: "concluye acerca del
%  cumplimiento/incumplimiento de los objetivos o intereses de la empresa")
%   >>> Vuelve a los TRES componentes del objetivo que declaraste en 1.3 y
%       pronúnciate sobre cada uno, uno por uno. El corrector va a buscar
%       exactamente esa correspondencia.
%
%   O1. Primer paso formal en capacidades espaciales -> CUMPLIDO.
%       Evidencia: existe un proyecto espacial formalizado, con documentación
%       normalizada bajo el sistema de identificación de documentos de DTS, y un
%       modelo de ingeniería que la empresa no poseía en enero de 2026.
%
%   O2. Alcanzar madurez de diseño PDR -> CUMPLIDO.
%       Evidencia: la PDR se cerró el 27 de febrero de 2026 con INVAP, dentro del
%       plazo de la práctica, con los entregables exigidos por la revisión.
%
%   O3. Aprovechar la transferencia tecnológica con INVAP -> CUMPLIDO PARCIALMENTE.
%       Evidencia: cuatro instancias formales de revisión con la contraparte, con
%       retroalimentación incorporada a los requerimientos y a la caracterización
%       del canal. Sin embargo, la transferencia en materia de fabricación,
%       integración y ensayos queda pendiente para las fases siguientes.
%       >>> Reconocer un cumplimiento parcial es más creíble que declarar éxito
%           total en todo, y la rúbrica admite "cumplimiento/INCUMPLIMIENTO".
%
% [ESTRUCTURA] 4 párrafos: uno introductorio que recuerda el objetivo y tres de
% pronunciamiento, uno por componente, cada uno con su evidencia.
%
% [EJEMPLO DE ARRANQUE]
%   "El objetivo que DTS definió para el período de práctica se cumplió en sus dos
%    primeros componentes y de forma parcial en el tercero. Respecto del primero,
%    la empresa cerró el período con un proyecto espacial formalizado..."

% ESCRIBIR AQUÍ


% ---------------------------------------------------------------------
\subsection{Impacto y relevancia de los resultados}
% ---------------------------------------------------------------------
% [EXTENSIÓN] ~200 palabras.
%
% [CONTENIDO]  (rúbrica B / Evaluación, criterio 4: justifica el impacto o la
%  relevancia -- eficiencia, productividad, clima laboral, etc.)
%   >>> Aquí hay que hablar de IMPACTO SOBRE LA EMPRESA, no de bondades del
%       satélite. Es el error más común en esta subsección.
%
%   Impacto en eficiencia y productividad de la ingeniería:
%     - El modelo MBSE sustituye la actualización manual de varios documentos por
%       una única fuente de verdad; el impacto se mide en horas de reconciliación
%       documental evitadas ante cada cambio de requerimiento.
%     - La matriz de trazabilidad permite estimar el impacto de un cambio antes
%       de ejecutarlo, en vez de descubrirlo en integración.
%
%   Impacto en reducción de riesgo económico:
%     - Cerrar la PDR antes de comprometer los 133.000 USD estimados evita
%       invertir en un diseño cuya viabilidad no estaba demostrada. El costo de
%       los dos meses de práctica es marginal frente a esa cifra.
%
%   Impacto en capacidad instalada:
%     - La empresa queda con tres activos reutilizables: el modelo Capella, la
%       herramienta de simulación de enlace en Python y una metodología de
%       revisiones documentada y replicable.
%     - Se formó capital humano en un área sin trayectoria previa, en línea con
%       el Plan Estratégico 2025-2029 de DTS.
%
%   Relevancia externa (una sola oración de cierre, no más):
%     - Aporte al ecosistema espacial nacional en línea con los objetivos del CEN.
%
% [ESTRUCTURA] 4 párrafos: eficiencia de la ingeniería, riesgo económico,
% capacidad instalada, relevancia externa.

% ESCRIBIR AQUÍ


% ---------------------------------------------------------------------
\subsection{Conclusiones}
% ---------------------------------------------------------------------
% [EXTENSIÓN] ~200 palabras. Máximo 4 párrafos.
%
% [CONTENIDO]
%   >>> NO introduzcas información nueva. Cierra la pregunta de la sección 1.3.
%   P1  Responde la pregunta planteada en 1.3: la ausencia de un proceso de
%       ingeniería preexistente se resolvió adoptando un marco normalizado de
%       revisiones (IEEE Std 15288.2) y una metodología de modelado (ARCADIA),
%       que aportaron simultáneamente la estructura del proceso y la trazabilidad
%       del producto.
%   P2  Resultado técnico principal: la arquitectura propuesta es viable. El
%       enlace cierra con margen en 4 a 5 pasadas diarias y la arquitectura basada
%       en SDR satisface los requerimientos críticos de radiación, modularidad y
%       actualización en vuelo.
%   P3  Lo que queda abierto: ensayos no ejecutados, escenario multi-DCP no
%       modelado, componentes por confirmar. La CDR es el siguiente umbral.
%   P4  Conclusión sobre el valor para la empresa: DTS pasó de no tener línea de
%       desarrollo espacial a contar con un proyecto formalizado, una metodología
%       replicable y herramientas propias reutilizables.
%
% [EJEMPLO DE ARRANQUE DE P1]
%   "El desarrollo de BIFROST permite responder la pregunta planteada al inicio de
%    este informe. La ausencia de un proceso de ingeniería de sistemas
%    preexistente en la empresa no se resolvió improvisando un método propio, sino
%    adoptando un marco normalizado de revisiones de madurez y una metodología de
%    modelado que impone trazabilidad por construcción."

% ESCRIBIR AQUÍ
